\section{Marco teórico y práctico}

El desarrollo de este TFG ha requerido la configuración y uso de un entorno de entrenamiento preexistente, al que se ha accedido mediante una librería, así como la modelización del problema objeto de estudio.

Para el entrenamiento del modelo y la evaluación de la complejidad del problema, resulta necesario caracterizarlo adecuadamente y modelarlo conforme al formato \emph{gen9randombattle}.

\subsection{Características generales}

La características generales del formato \emph{gen9randombattle} que son fundamentales para la modelización del problema mediante aprendizaje por refuerzo son las siguientes:

\begin{itemize}

    \item Decisiones secuenciales

          Cada batalla se desarrolla por turnos discretos en los que hay que elegir una acción. Las acciones tomadas afectan al estado de la batalla, lo que hace que sea secuencial.

    \item Adversarial

          En cada batalla hay dos jugadores que compiten y solo puede ganar uno de ellos. Como no hay cooperación, el problema es de suma cero.

    \item Información oculta

          El agente dispone de información completa sobre su propio equipo, mientras que la información del oponente es parcial y se va revelando progresivamente durante el combate, a partir de las observaciones y el historial de acciones. Inicialmente, únicamente se observa el Pokémon de salida del rival y los posibles efectos activados al entrar en el campo. En la mayoría de las ocasiones el combate termina sin que se llegue a conocer la composición completa del equipo contrario.

    \item Alto grado de estocasticidad

          Una vez elegidas las acciones, el resultado del turno no es determinista. Cada ataque tiene una probabilidad de acertar, una de golpe crítico (alrededor de 1/24), algunos tienen posibles efectos secundarios y un rango aleatorizado de daño (se multiplica por un entero aleatorio entre 85 y 100 y se divide por 100). Por turno puede haber dos ataques y además hay ciertas habilidades o objetos cuya activación es aleatoria. Todo esto hace que pueda haber turnos en los cuales el resultado es completamente predecible, meintras que en otros el resultado presenta alto grado de estocasticidad.

    \item Acciones simultáneas

          Al inicio de cada turno, ambos jugadores eligen su acción sin conocer la acción del oponente, lo que lleva a situaciones en las que conseguir la victoria dependa de predecir correctamente la acción del oponente. Un ejemplo de un juego que responde a este patrón es el  famoso piedra-papel-tijera.

    \item Recompensas retrasadas

          La recompensa se recibe al final de la batalla únicamente cuando todos los Pokémon de un jugador han sido derrotados. Se podría pensar en intentar medir los recursos restantes o el estado de la batalla para dar recompensas intermedias, pero debido a la alta complejidad del entorno, esto no se puede realizar de manera objetiva o cuantificable.

\end{itemize}

\subsection{Marco teórico del problema}

En esta sección se describe el problema de las batallas del formato \emph{gen9randombattle} como POMDPs (\textit{Partially Observable Markov Decision Process}), o Proceso de decisión de Markov parcialmente observable. Para ello hay que definir el estado, las acciones y las transiciones. La tabla \ref{tab:caracteristicas} contiene la nomenclatura que se utilizará en el resto de esta sección. Si se desea ver una modelización más completa y detallada, se puede consultar el anexo \ref{anexo:modeling}.

\begin{table}[h]
    \centering
    \begin{tabular}{ll}
        \toprule
        Símbolo                 & Significado                                                                      \\
        \midrule
        $t \in \{0,1,2,\dots\}$ & Índice del turno.                                                                \\
        $\cS$                   & Espacio de estados completo oculto de la batalla.                                \\
        $s_t \in \cS$           & Estado markoviano completo en el turno $t$.                                      \\
        $i \in \{1,2\}$         & Índice del jugador.                                                              \\
        $\cA_i(s)$              & Conjunto de acciones legales para el jugador $i$ en el estado $s$.               \\
        $a_t^i \in \cA_i(s_t)$  & Acción elegida por el jugador $i$ en el turno $t$.                               \\
        $\gamma \in (0,1]$      & Factor de descuento.                                                             \\
        $\Xi$                   & Conjunto de etiquetas de especies o tipos.                                       \\
        $\cM$                   & Espacio de configuraciones latentes de movimientos.                              \\
        $\cU$                   & Espacio de configuraciones latentes de objetos/habilidades/estadísticas.         \\
        $\mathsf{HP}$           & Variables de puntos de vida para todas las unidades activas y de reserva.        \\
        $\mathsf{Field}$        & Estado público del campo (clima, terreno, trampas, pantallas, contadores, etc.). \\
        \bottomrule
    \end{tabular}
    \caption{Terminología y formulación usados en el modeloado de un combate Pokémon  \emph{gen9randombattle} como un POMDP}
    \label{tab:caracteristicas}
\end{table}

\subsubsection{Descomposición del estado}

El estado completo de la batalla debe ser lo suficientemente rico para que el juego sea markoviano. Por tanto, definimos
$$
    \cS = \cS_{\mathrm{team},1} \times \cS_{\mathrm{team},2} \times \cS_{\mathrm{battle}}
$$
donde $\cS_{\mathrm{team},i}$ almacena la información a nivel de equipo del jugador $i$, y $\cS_{\mathrm{battle}}$ guarda la configuración de la batalla en evolución.

Una factorización conveniente es
$$
    s_t = \big(T_1^t,T_2^t,B_t\big)
$$
donde:
\begin{itemize}
    \item $T_i^t$ es el estado completo del equipo del jugador $i$ en el turno $t$.
    \item $B_t$ es el estado compartido de la batalla (campo, clima, trampas, contadores de turno, efectos pendientes, etc.).
\end{itemize}

Cada estado de equipo puede escribirse como
$$
    T_i^t = \big( X_{i,1}^t, \dots, X_{i,6}^t, A_i^t \big)
$$
donde $X_{i,k}^t$ es el estado detallado de la $k$-ésima unidad, y $A_i^t$ indica qué unidad está actualmente activa.

Cada estado de unidad tiene componentes cuasi-estáticos y dinámicos:
$$
    X_{i,k}^t = \big(\xi_{i,k}, m_{i,k}, u_{i,k}, y_{i,k}^t\big)
$$
donde
\begin{itemize}
    \item $\xi_{i,k} \in \Xi$ es la etiqueta de especie o arquetipo.
    \item $m_{i,k} \in \cM$ denota la especificación del conjunto de movimientos.
    \item $u_{i,k} \in \cU$ denota objetos latentes, habilidad, naturaleza, parámetros de estadísticas u otra configuración oculta.
    \item $y_{i,k}^t$ denota la condición interna dependiente del tiempo de la unidad.
\end{itemize}

Una variable de condición interna detallada puede escribirse como
$$
    y_{i,k}^t = \big(\mathrm{hp}_{i,k}^t,\mathrm{status}_{i,k}^t,\mathrm{boosts}_{i,k}^t,\mathrm{volatile}_{i,k}^t,\mathrm{fainted}_{i,k}^t\big)
$$

El estado compartido de la batalla puede igualmente escribirse como
$$
    B_t = \big(\mathrm{weather}_t,\mathrm{terrain}_t,\mathrm{hazards}_t,\mathrm{screens}_t,\mathrm{trickroom}_t,\mathrm{timers}_t,\dots\big)
$$

\subsubsection{Acciones y restricciones legales}

En el turno $t$, ambos jugador eligen una acción simultáneamente y a priori. Para el jugador $i$, el conjunto de acciones legales depende del estado completo actual:
$$
    a_t^i \in \cA_i(s_t)
$$

En \emph{gen9randombattle}, un conjunto de acciones legales típico tiene la forma esquemática
$$
    \cA_i(s_t) \subseteq \cA_i^{\mathrm{move}}(s_t) \cup \cA_i^{\mathrm{switch}}(s_t) \cup \cA_i^{\mathrm{tera}}(s_t)
$$

donde las acciones de movimiento corresponden a seleccionar un movimiento legal, las acciones de cambio corresponden a cambiar a una unidad de reserva legal, y donde $\cA_i^{\mathrm{tera}}(s_t) = \{0,1\}$ denota la decisión booleana de si Terascristalizar o no. Esto solo se puede hacer una vez por batalla por jugador y debe elegirse con un movimiento y no con un cambio.

\subsubsection{Transiciones}

El motor del juego define un núcleo de transición estocástico
\begin{equation}
    P(s_{t+1}\mid s_t,a_t^1,a_t^2)
    \label{eq:transition_kernel}
\end{equation}
Este núcleo incluye toda la resolución basada en reglas, incluyendo el orden por prioridad y velocidad, la ejecución de movimientos, cambios, precisión, variación de daño, golpes críticos, efectos secundarios, activaciones de efectos, actualizaciones al final del turno y debilitamientos o cambios forzados cuando corresponda.

El siguiente estado evoluciona por tanto como
$$
    s_{t+1} \sim P(\cdot\mid s_t,a_t^1,a_t^2)
$$

Nota sobre acciones simultaneas y resolución interna: aunque los jugadores eligen sus acciones de manera simultánea a nivel estratégico, el motor del juego resuelve la acción conjunta siguiendo un orden interno que puede ser complejo; dicho orden forma parte del núcleo de transición y no necesita representarse por separado una vez definida la Ecuación~\eqref{eq:transition_kernel}.

\subsection{Entorno de entrenamiento}

Para jugar múltiples partidas se necesita un simulador. Aunque no es oficial de Game Freak (la compañía que desarrolla Pokémon, bajo el control de Nintendo), el más popular es Pokémon Showdown~\cite{pokemonShowdown}. Pokémon Showdown es un simulador de código abierto que no tiene una API oficial, pero existe la librería de python poke-env~\cite{poke_env} que permite conectarse a un servidor de Pokémon Showdown y jugar partidas. El servidor no permite repetir partidas (estocástica no repetible) ni empezar desde un punto intermedio, lo que hace el problema más difícil.


Dado que el entrenamiento requiere la ejecución de millones de partidas, no es viable utilizar el servidor público de Pokémon Showdown. No obstante, el proyecto se distribuye bajo licencia \textit{MIT license}, lo que permite desplegar un servidor local y privado sin restricciones de uso ni de comunicación. En este trabajo se ha empleado la última versión estable disponible de Pokémon Showdown, correspondiente a la versión v0.11.10. No se ha utilizado la versión directamente disponible en GitHub, ya que esta se actualiza de forma continua, lo que impide garantizar su estabilidad.

La librería poke-env viene preparada para conectarse a cualquier servidor y permitir que un agente juegue múltiples partidas a la vez. A continuación se incluye el código para conectar a un servidor:

\begin{verbatim}
ServerConfiguration(
    f"ws://localhost:{int(os.getenv('SERVER_PORT'))}/showdown/websocket",
    "https://play.pokemonshowdown.com/action.php?",
)
\end{verbatim}

Para que el agente pueda jugar, simplemente tiene que heredar de la clase \texttt{poke\_env.player.Player} e implementar el método \texttt{choose\_move}. Es en este método donde se recibe el estado de la batalla. Este estado de la batalla será la entrada al modelo (la red neuronal), que devolverá la acción a realizar que a su vez será la salida del método \texttt{choose\_move}.
